\newpage


\section{SYS\_06 - Database security}


\subsection{Introduzione alla sicurezza dei database}

La sicurezza dei database costituisce l’insieme di meccanismi, controlli e pratiche operative finalizzate alla protezione dei dati conservati all’interno di un DBMS.
L’obiettivo centrale è difendere il patrimonio informativo da minacce intenzionali (attacchi informatici, accessi non autorizzati, insider threats) e da eventi accidentali (errori umani, malfunzionamenti, perdita di disponibilità del sistema) \footnote{Vedi ``What is database security?'' nel PDF.}.

Una \textit{security policy} definisce in modo rigoroso le misure da adottare, mentre il DBMS fornisce gli strumenti per farla rispettare.
Poiché la superficie d'attacco si estende oltre il DBMS, risulta necessario integrare tali meccanismi con procedure sistemistiche, di rete e gestionali, così da proteggere l'intero stack che ospita i dati.

\subsection{Obiettivi e ambito della sicurezza}

Il documento enfatizza come il tema della sicurezza nei corsi classici di database sia spesso trattato marginalmente. Tuttavia, nei sistemi reali diventa un’area critica, che coinvolge:

\begin{itemize}
    \item l’insieme dei dati memorizzati nel DB;
    \item il DBMS come software (comprensivo delle sue vulnerabilità);
    \item le applicazioni che lo utilizzano;
    \item il server fisico o virtuale in cui risiede il DBMS;
    \item l'infrastruttura di rete utilizzata per accedere ai dati.
\end{itemize}

Un concetto chiave presentato nel PDF è il \textbf{paradosso di Anderson}: aumentando l’accessibilità del database se ne riduce la sicurezza; restringendo e irrobustendo le difese, si può compromettere l’usabilità.

\subsection{Minacce comuni}

Le principali minacce individuate includono \footnote{``Common threats'' nel PDF.}:

\begin{itemize}
    \item \textbf{Insider threats}, dovute a utenti legittimi che abusano dei propri privilegi.
    \item \textbf{Errori umani}, come configurazioni errate, query mal formulate o cancellazioni involontarie.
    \item \textbf{Vulnerabilità del DBMS}, sfruttabili da attaccanti esterni.
    \item \textbf{Vulnerabilità applicative}, in particolare SQL injection.
    \item \textbf{Buffer overflow} e altre categorie di exploit.
    \item \textbf{Malware} che compromette i sistemi host.
    \item \textbf{Attacchi ai backup}, se non adeguatamente protetti.
    \item \textbf{DDoS}, che colpiscono disponibilità e prestazioni.
\end{itemize}

\subsection{Best Practices generali}

Le best practices della sicurezza database includono \footnote{``Best practices'' nel PDF.}:

\begin{itemize}
    \item sicurezza fisica delle infrastrutture;
    \item controllo degli accessi (PoLP);
    \item sicurezza degli account utente e dei dispositivi;
    \item cifratura dei dati e dei canali di comunicazione;
    \item hardening del DBMS;
    \item messa in sicurezza dell'applicazione e del web server;
    \item protezione dei backup;
    \item auditing e monitoraggio.
\end{itemize}

\subsection{Sicurezza fisica}

\subsubsection{Difesa a più livelli}

Le slide mostrano chiaramente come i data center applichino un modello multilivello, comprendente perimetro esterno, guardiania, accesso agli edifici, zone sicure, controllo ascensori, corridoi protetti e accesso ai rack.
Un aspetto cruciale è la gestione dei supporti di memorizzazione: gli HDD vengono tracciati e distrutti fisicamente a fine vita, come illustrato nella ``destruction room'' di Google \footnote{``Physical security'' pp. 9--10.}.

\subsection{SQL Injection}

\subsubsection{Panoramica}

Le SQL injection rappresentano una delle minacce più critiche ai sistemi basati su DBMS.
Esse derivano dall’inserimento non sicuro di input dell’utente all’interno di query.
Un attaccante può alterare la semantica della query, accedere a dati sensibili, modificare tabelle o compromettere il server.

\subsubsection{Tipologie}

\begin{itemize}
    \item \textbf{First-order SQLi}: l'input dell’utente viene immediatamente inserito nella query vulnerabile.
    \item \textbf{Second-order SQLi}: l'input pericoloso è salvato nel database e sfruttato successivamente, quando recuperato dall’applicazione.
\end{itemize}

\subsubsection{Tecniche di rilevazione}

Le tecniche includono l’invio sistematico di payload quali:

\begin{itemize}
    \item apici per rompere la struttura sintattica (\verb+'+);
    \item condizioni booleane (\verb+OR 1=1+);
    \item payload basati su ritardi temporali;
    \item interazioni out-of-band (OAST).
\end{itemize}

Il PDF illustra numerosi esempi pratici (bypass login, esposizione dati nascosti, estrazione tramite \texttt{UNION}) e link a laboratori PortSwigger.

\subsubsection{Prepared Statements}

Il documento sottolinea l’importanza dei \textit{prepared statements} come tecnica per prevenire SQLi: il DBMS precompila il piano di esecuzione della query con placeholder, separando codice e dati.

\subsection{Autenticazione, autorizzazione e controllo degli accessi}

\subsubsection{Autenticazione}

L'autenticazione stabilisce se l’utente è realmente chi afferma di essere.
I DBMS moderni richiedono che solo utenti autenticati possano creare connessioni.

\subsubsection{Autorizzazione}

L'autorizzazione assegna i privilegi che determinano quali operazioni un utente può svolgere.
Il modello richiamato è prevalentemente \textbf{DAC}, dove il creatore della risorsa detiene diritti completi.

\subsubsection{Privilegi SQL}

I principali privilegi includono:

\begin{itemize}
    \item \texttt{SELECT}, \texttt{INSERT}, \texttt{UPDATE}, \texttt{DELETE};
    \item \texttt{REFERENCES} per definire chiavi esterne;
    \item \texttt{USAGE} su domini;
    \item \texttt{EXECUTE} su stored procedures.
\end{itemize}

Vengono inoltre trattati i comandi \texttt{GRANT}, \texttt{REVOKE} e l’importanza dei ruoli introdotti da SQL-99.

\subsubsection{Viste come meccanismo di sicurezza}

Le viste consentono di esporre solo una porzione del database, implementando un controllo fine (column-level o row-level), senza che l’utente veda ciò che non è autorizzato.

\subsection{RAID e affidabilità dello storage}

La sicurezza non è solo confidenzialità: include disponibilità e resilienza.
La sezione dedicata al RAID analizza livelli da 0 a 6 e RAID-10, illustrando compromessi tra throughput, concorrenza di accesso e tolleranza ai guasti.

\subsubsection{Punti chiave}

\begin{itemize}
    \item RAID-0: massime prestazioni, nessuna ridondanza.
    \item RAID-1: mirroring, ottima tolleranza ai guasti.
    \item RAID-5: parità distribuita, buon compromesso.
    \item RAID-6: doppia parità, resiste a due guasti.
    \item RAID-10: ideale per RDBMS con molti accessi casuali.
\end{itemize}

\subsection{Backup e Disaster Recovery}

Il PDF enfatizza che i backup devono essere soggetti a precise policy: posizione, accesso, cifratura, rotazione.

\subsubsection{Disaster Recovery}

Vengono presentati due indicatori fondamentali:

\begin{itemize}
    \item \textbf{RPO} (Recovery Point Objective): quantità di dati che si può perdere.
    \item \textbf{RTO} (Recovery Time Objective): tempo necessario a ripristinare l’operatività.
\end{itemize}

È ricordata la \textbf{regola 3-2-1}: tre copie dei dati, due supporti diversi, una off-site.

\subsection{Cifratura e Homomorphic Encryption}

La cifratura dei dati a riposo e in transito è fondamentale.
Il documento estende poi la discussione all’\textbf{homomorphic encryption}, che consente il calcolo su dati cifrati senza decrittarli.

\subsubsection{Fully Homomorphic Encryption}

Viene introdotta la distinzione tra:

\begin{itemize}
    \item \textbf{Somewhat homomorphic encryption}: supporta solo un numero limitato di operazioni prima che il rumore diventi eccessivo;
    \item \textbf{Leveled FHE}: organizzata in livelli con bootstrapping per eliminare il rumore.
\end{itemize}

\subsection{Database Monitoring}

Il monitoraggio copre tre aree:

\begin{itemize}
    \item \textbf{Performance}: tempi di risposta, uso risorse, errori.
    \item \textbf{Security}: accessi anomali, modifiche sospette.
    \item \textbf{Compliance}: conformità a standard quali GDPR e HIPAA.
\end{itemize}

Le slide citano l’uso di \texttt{USDT/uprobes} ed \textbf{eBPF} per osservabilità avanzata del comportamento del DBMS.

\subsection{DBMS Security}

Il documento conclude ricordando che il DBMS è software, e come tale può contenere vulnerabilità: è quindi necessario applicare patch e aggiornamenti in modo tempestivo.
